\section{The adelic big picture}\label{sec:bigpicture}

Conway's big picture \cite{Conway1996} is the set of homothety classes of
lattices in $\QQ^{2}$ with its natural action of $\RationalProjGroup$. We
describe it adelically.

\subsection{Lattices and vertices}

\begin{definition}[Lattices]\label{def:lattices}
A \emph{lattice} in $\StdModule{\FiniteAdeles{\QQ}}{2}$ is a compact open
$\ProfiniteIntegers$-submodule. For a prime $p$ a lattice in
$\StdModule{\pAdicNumbers{p}}{2}$ is a compact open
$\pAdicIntegers{p}$-submodule.
\end{definition}

\begin{lemma}[Lattices are translates of the standard lattice]
\label{lem:lattices-translates}
The lattices in $\StdModule{\FiniteAdeles{\QQ}}{2}$ are the sets
$g\StdModule{\ProfiniteIntegers}{2}$ with $g \in \GL{2}{\FiniteAdeles{\QQ}}$, the
stabiliser of $\StdModule{\ProfiniteIntegers}{2}$ is
$\GL{2}{\ProfiniteIntegers}$, and $L \mapsto (L\otimes\pAdicIntegers{p})_p$ is
a bijection onto the families of lattices $L_p \subseteq
\StdModule{\pAdicNumbers{p}}{2}$ with $L_p = \StdModule{\pAdicIntegers{p}}{2}$
for almost all $p$. The map $L \mapsto L \otimes \ProfiniteIntegers$ is a
bijection from lattices in $\StdModule{\QQ}{2}$ to lattices in
$\StdModule{\FiniteAdeles{\QQ}}{2}$.
\end{lemma}

\begin{proof}
Standard; the last statement is class number one for $\GL{2}{}$
(\cref{lem:gl2-class-number-one}).
\end{proof}

\begin{definition}[Vertices of the big picture]\label{def:big-picture-vertices}
Let $\BigPictureVertices$ be the set of lattices in
$\StdModule{\FiniteAdeles{\QQ}}{2}$ modulo homothety by
$\unit{\FiniteAdeles{\QQ}}$, and write $\LatticeClass{L}$ for the class of $L$.
For a prime $p$ let $\LocalVertexSet{p}$ be the set of lattices in
$\StdModule{\pAdicNumbers{p}}{2}$ modulo $\unit{\pAdicNumbers{p}}$; it is the
vertex set of the Bruhat--Tits tree $\Tree{p}$ of $\PGL{2}{\pAdicNumbers{p}}$,
a $(p+1)$-regular tree in which two classes are adjacent if they have
representatives $L \supsetneq L' \supsetneq pL$. The \emph{base vertex} is
$\BaseVertex := \LatticeClass{\StdModule{\ProfiniteIntegers}{2}}$ and
$\LocalBaseVertex{p} := \LatticeClass{\StdModule{\pAdicIntegers{p}}{2}}$.
For $t \in \PositiveRationals$ let
$\StdLattice{t} := \ProfiniteIntegers \oplus t\ProfiniteIntegers =
\Diagonal{1}{t}\StdModule{\ProfiniteIntegers}{2}$ (column vectors) and
$\StdVertex{t} := \LatticeClass{\StdLattice{t}}$, with local components
$\LocalStdVertex{t}{p}$.
\end{definition}

\begin{lemma}[The big picture is a restricted product of trees]
\label{lem:restricted-product}
The map $\LatticeClass{L} \mapsto (\LatticeClass{L\otimes\pAdicIntegers{p}})_p$
is a bijection from $\BigPictureVertices$ onto the restricted product of
the $\LocalVertexSet{p}$ with respect to the $\LocalBaseVertex{p}$. Via
\cref{lem:lattices-translates} it identifies $\BigPictureVertices$ with the
set of lattices in $\QQ^{2}$ modulo $\unit{\QQ}$, i.e.\ with Conway's big
picture (in the column convention).
\end{lemma}

\begin{proof}
The first statement follows from \cref{lem:lattices-translates} because
$\unit{\FiniteAdeles{\QQ}}$ is the restricted product of the
$\unit{\pAdicNumbers{p}}$ with respect to the $\unit{\pAdicIntegers{p}}$. For
the second, $\unit{\FiniteAdeles{\QQ}} = \PositiveRationals\unit{\ProfiniteIntegers}$
(\cref{lem:rational-scalars-discrete}) and $\unit{\ProfiniteIntegers}$ fixes
every lattice.
\end{proof}

\begin{definition}[Hyperdistance]\label{def:hyperdistance}
For $x, y \in \BigPictureVertices$ let
$\HyperDistance{x}{y} := \prod_p p^{\TreeDistance{p}{x_p}{y_p}}$, where
$\TreeDistance{p}{}{}$ is the distance in $\Tree{p}$.
\end{definition}

\begin{lemma}[Hyperdistance as an index]\label{lem:hyperdistance-index}
Let $x, y \in \BigPictureVertices$. There are representatives $L \supseteq L'$
of $x$ and $y$ with $L/L'$ cyclic, and then
$\HyperDistance{x}{y} = \#(L/L')$. For $t \mid t'$ positive integers,
$\HyperDistance{\StdVertex{t}}{\StdVertex{t'}} = t'/t$.
\end{lemma}

\begin{proof}
This is the corresponding standard statement for $\Tree{p}$, applied to
each prime separately.
\end{proof}

\begin{definition}[Geometric realisation]\label{def:big-picture}
Let $\TreeRealization{p}$ be the geometric realisation of $\Tree{p}$, in
which every edge is isometric to $[0,1]$. The \emph{big picture} is
\[
  \BigPicture := \Bigl\{ (x_p)_p \in \prod_p \TreeRealization{p} :
  x_p = \LocalBaseVertex{p} \text{ for almost all } p \Bigr\} .
\]
The group $\AdelicGroup$ acts on $\BigPicture$ through
$\ProjectionToPGL \colon \AdelicGroup \to \PGL{2}{\FiniteAdeles{\QQ}}$, the
component at $p$ acting on $\TreeRealization{p}$ by isometries. For
$x \in \TreeRealization{p}$ and adjacent vertices $u, v$ write
$\Midpoint{u}{v}$ for the midpoint of the edge between them.
\end{definition}

\begin{lemma}[Stabiliser of the base vertex]\label{lem:stab-base-vertex}
\slogan{The stabiliser of the standard lattice class is
$\GL{2}{\ProfiniteIntegers}$.}
$\Stab_{\AdelicGroup}(\BaseVertex) = \AdelicProjection(\GL{2}{\ProfiniteIntegers})$.
\end{lemma}

\begin{proof}
If $g\StdModule{\ProfiniteIntegers}{2} = \lambda\StdModule{\ProfiniteIntegers}{2}$
with $\lambda \in \unit{\FiniteAdeles{\QQ}}$, write $\lambda = qu$ with
$q \in \PositiveRationals$ and $u \in \unit{\ProfiniteIntegers}$
(\cref{lem:rational-scalars-discrete}); then $q^{-1}g \in \GL{2}{\ProfiniteIntegers}$.
\end{proof}

\subsection{The standard groups}

\begin{definition}[Adelic standard groups]\label{def:adelic-standard-groups}
Let $N \geq 1$. Define the following subgroups of $\AdelicGroup$, all
contained in $\AdelicProjection(\GL{2}{\ProfiniteIntegers})$:
\begin{align*}
  \GammaZero[ad]{N} &:= \AdelicProjection\{\smallMat{a}{b}{c}{d} \in \GL{2}{\ProfiniteIntegers} : c \in N\ProfiniteIntegers\},\\
  \GammaOne[ad]{N} &:= \AdelicProjection\{g \in \GL{2}{\ProfiniteIntegers} : g \equiv \smallMat{1}{*}{0}{*} \bmod N\},\\
  \GammaFull[ad]{N} &:= \AdelicProjection\{g \in \GL{2}{\ProfiniteIntegers} : g \equiv 1 \bmod N\}.
\end{align*}
Their lifted and classical versions (\cref{def:versions}) are denoted
$\GammaZero[lift]{N}$, $\GammaZero{N}$, and so on.
\end{definition}

\begin{lemma}[Classical versions of the standard groups]
\label{lem:classical-standard-groups}
Let $N \geq 1$. Then $\GammaZero[lift]{N}$, $\GammaOne[lift]{N}$ and
$\GammaFull[lift]{N}$ are the usual subgroups of $\SL{2}{\ZZ}$ defined by
$c \equiv 0$, by $c \equiv 0, a \equiv d \equiv 1$, and by
$\gamma \equiv 1$ modulo $N$, and $\GammaZero{N}$, $\GammaOne{N}$,
$\GammaFull{N}$ are their images in $\PSL{2}{\ZZ} \subseteq \RationalProjGroup$.
\end{lemma}

\begin{proof}
By the description of rational elements of $\GL{2}{\ProfiniteIntegers}$
(\cref{lem:rational-points-glzhat}); for $\GammaOne{}$ note that $a d \equiv 1$
when $c \equiv 0$ and $\det = 1$.
\end{proof}

\begin{lemma}[{$\GammaZero[ad]{N}$ fixes the snake}]\label{lem:gamma0-fixer}
\slogan{$\GammaZero{N}$ is the fixer of the lattice classes
$\StdVertex{1}$ and $\StdVertex{N}$.}
Let $N \geq 1$ and $t$ a positive integer. An element
$g = \smallMat{a}{b}{c}{d} \in \GL{2}{\ProfiniteIntegers}$ fixes
$\StdVertex{t}$ if and only if $c \in t\ProfiniteIntegers$. Consequently
\[
  \GammaZero[ad]{N} = \Stab_{\AdelicGroup}(\StdVertex{1}) \cap \Stab_{\AdelicGroup}(\StdVertex{N}),
\]
and $\GammaZero[ad]{N}$ fixes $\StdVertex{d}$ for every divisor $d$ of $N$.
\end{lemma}

\begin{proof}
One has $g\smallColVec{1}{0} = \smallColVec{a}{c}$ and
$g\smallColVec{0}{t} = \smallColVec{tb}{td}$, so
$g\StdLattice{t} \subseteq \StdLattice{t}$ if and only if
$c \in t\ProfiniteIntegers$, and then equality holds since $g^{-1}$ has the
same shape. Combine with the stabiliser of the base vertex
(\cref{lem:stab-base-vertex}).
\end{proof}

\begin{definition}[Snake]\label{def:snake}
For $N \geq 1$ the \emph{snake} of $N$ is
$\Snake{N} := \{\StdVertex{d} : d \mid N\} \subseteq \BigPictureVertices$.
Its local component at $p$ is the geodesic segment from
$\LocalBaseVertex{p}$ to $\LocalStdVertex{N}{p}$, of length $\val{p}{N}$.
\end{definition}

\begin{definition}[Local flips and the Atkin--Lehner group]
\label{def:atkin-lehner-group}
Let $N \geq 1$ and let $p$ be a prime with $k := \val{p}{N} \geq 1$. The
\emph{local flip} $\ALflip{p,N} \in \AdelicGroup$ is the image of the
element of $\GL{2}{\FiniteAdeles{\QQ}}$ equal to
$\smallMat{0}{-1}{p^{k}}{0}$ at $p$ and to $1$ at all other primes. The
\emph{Atkin--Lehner group} is the setwise stabiliser
\[
  \GammaZeroStar[ad]{N} := \Stab_{\AdelicGroup}(\Snake{N}).
\]
\end{definition}

\begin{lemma}[Structure of the Atkin--Lehner group]
\label{lem:atkin-lehner-group}
\slogan{The setwise stabiliser of the snake is $\GammaZero{N}$ extended by
independent flips at the primes dividing $N$.}
Let $N \geq 1$. The local flip $\ALflip{p,N}$ preserves the local snake at
$p$ and reverses it. The pointwise stabiliser of $\Snake{N}$ is
$\GammaZero[ad]{N}$, and the map recording at which primes an element
reverses the snake induces an isomorphism
$\GammaZeroStar[ad]{N}/\GammaZero[ad]{N} \cong \bigoplus_{p \mid N}\ZZ/2\ZZ$,
with $\ALflip{p,N}$ mapping to the generator at $p$. If $N$ is squarefree,
this isomorphism is the type map
$\GammaZeroStar[ad]{N} \to \TypeGroup{N}$.
\end{lemma}

\begin{proof}
At $p$ one computes $\smallMat{0}{-1}{p^{k}}{0}\smallColVec{x}{y} =
\smallColVec{-y}{p^{k}x}$, so the flip maps
$\StdModule{\pAdicIntegers{p}}{2}$ to $\pAdicIntegers{p}\oplus p^{k}\pAdicIntegers{p}$
and $\pAdicIntegers{p}\oplus p^{j}\pAdicIntegers{p}$ to the class of
$\pAdicIntegers{p}\oplus p^{k-j}\pAdicIntegers{p}$. An isometry
preserving a segment either fixes it pointwise or reverses it, and it fixes
it pointwise if it fixes one endpoint. The pointwise stabiliser is
$\GammaZero[ad]{N}$ by \cref{lem:gamma0-fixer}. For squarefree $N$ the
element $\ALflip{p,N}$ has determinant of valuation $1$ at $p$, and elements
of $\GammaZero[ad]{N}$ have trivial type.
\end{proof}

\begin{definition}[Global Atkin--Lehner matrices]\label{def:global-al}
Let $N \geq 1$ and let $e$ be a Hall divisor of $N$, i.e.\ $e \mid N$ and
$\gcd(e, N/e) = 1$. An \emph{Atkin--Lehner matrix} for $e$ is a matrix
$\ALmatrix{e} = \smallMat{ex}{y}{Nz}{ew}$ with $x,y,z,w \in \ZZ$ and
$\det\ALmatrix{e} = e$.
\end{definition}

\begin{lemma}[Atkin--Lehner matrices exist and flip]
\label{lem:global-al}
\slogan{The lifted Atkin--Lehner group is generated by $\GammaZero{N}$ and
the classical Atkin--Lehner matrices.}
Let $N \geq 1$ and $e$ a Hall divisor of $N$. An Atkin--Lehner matrix
$\ALmatrix{e}$ exists, and $\AdelicProjection(\ALmatrix{e}) \in
\GammaZero[ad]{N}\prod_{p \mid e}\ALflip{p,N}$. Moreover
\[
  \GammaZeroStar[lift]{N} = \bigcup_{e} \GammaZero[lift]{N}\,\ALmatrix{e},
\]
the union over the Hall divisors $e$ of $N$, for any choice of the
$\ALmatrix{e}$.
\end{lemma}

\begin{proof}
Since $\gcd(e, N/e) = 1$ choose $w$ with $ew \equiv 1 \bmod N/e$ and put
$x = 1$, $z = 1$, $y = (ew-1)/(N/e)$; then $e^{2}xw - Nyz = e$. Locally at
$p \mid e$ with $k = \val{p}{N}$ one computes
$\ALmatrix{e}\smallMat{0}{-1}{p^{k}}{0}^{-1} =
\smallMat{-y}{ex p^{-k}}{-ew}{Nzp^{-k}}$, which is integral with lower left
entry divisible by $p^{k}$ and unit determinant $e/p^{k}$. At primes
$p \nmid e$ the matrix $\ALmatrix{e}$ lies in $\GL{2}{\pAdicIntegers{p}}$
with lower left entry divisible by $p^{\val{p}{N}}$. This proves the first
two statements. If $g \in \GammaZeroStar[lift]{N}$ reverses the snake
exactly at the primes dividing $e$, then $g\ALmatrix{e}^{-1}$ is a rational
element of $\GammaZero[ad]{N}$, i.e.\ an element of $\GammaZero[lift]{N}$ by
\cref{lem:classical-standard-groups}.
\end{proof}

\begin{remark}
The classical group $\GammaZeroStar{N}$ is the group usually denoted
$\GammaZero{N}+$ (all Atkin--Lehner involutions), in Conway's notation
\cite{Conway1996}. Its subgroups containing $\GammaZero{N}$ are the groups
$\GammaZero{N}+e,f,\ldots$.
\end{remark}

\begin{definition}[Conway's groups]\label{def:conway-groups}
Let $h, n$ be positive integers with $h \mid n$. Define
\[
  \GammaConway[ad]{n}{h} := \Stab_{\AdelicGroup}(\StdVertex{h}) \cap
  \Stab_{\AdelicGroup}(\StdVertex{n}).
\]
\end{definition}

\begin{lemma}[Conway's groups are conjugates of $\GammaZero{}$]
\label{lem:conway-groups}
Let $h \mid n$ be positive integers and $D := \Diagonal{1}{h}$. Then
$\GammaConway[ad]{n}{h} = D\,\GammaZero[ad]{n/h}\,D^{-1}$, and
\[
  \GammaConway[lift]{n}{h} = \Bigl\{\smallMat{a}{b/h}{nc}{d} :
  a,b,c,d \in \ZZ,\ ad - \tfrac{n}{h}bc = 1\Bigr\}.
\]
\end{lemma}

\begin{proof}
$D\Diagonal{1}{t} = \Diagonal{1}{ht}$, so $D$ maps $\StdVertex{1}$ and
$\StdVertex{n/h}$ to $\StdVertex{h}$ and $\StdVertex{n}$; apply
$\GammaZero{}$ as a fixer (\cref{lem:gamma0-fixer}). Since $D$ is rational,
$\GammaConway[lift]{n}{h} = D\GammaZero[lift]{n/h}D^{-1}$, and
$D\smallMat{a}{b}{(n/h)c}{d}D^{-1} = \smallMat{a}{b/h}{nc}{d}$.
\end{proof}

\begin{remark}[Conway's normaliser theorem]\label{rem:conway-normaliser}
Conway \cite{Conway1996} shows that the normaliser of $\GammaZero{N}$ in
$\PSL{2}{\RR}$ is $\GammaConway{N/h}{h}+$, the extension of
$\GammaConway{N/h}{h}$ by all its Atkin--Lehner involutions, where $h$ is
the largest divisor of $24$ with $h^{2} \mid N$ (see also
\cite{AtkinLehner1970}). The number $24$ enters because $h \mid 24$ is
equivalent to $u^{2} \equiv 1 \bmod h$ for all units $u$. We do not reprove
this here.
\todo[inline]{Conway works with row vectors; check that his labelling of
the fixed classes of $\GammaConway{n}{h}$ agrees with
\cref{def:conway-groups}.}
\end{remark}
